% zakljucak.tex - ZAKLJUČAK

\section{ZAKLJUČAK}
\label{sec:zakljucak}

U ovom seminarskom radu detaljno su istraženi i analizirani upravljački pristupi za autonomno robotsko ultrazvučno snimanje, s posebnim naglaskom na kritičnu ulogu regulacije kontaktne sile između ultrazvučne sonde i pacijentova mekog tkiva. Kroz sustavni pregled literature i teorijsku analizu, pokazano je da pozicijsko upravljanje nije prikladno za medicinsko okruženje zbog rizika za sigurnost i nemogućnosti kompenzacije dinamike tijela.

Nasuprot tome, neizravno upravljanje silom -- prvenstveno upravljanje impedancijom i admitancijom -- te hibridno pozicijsko/silno upravljanje čine osnovu modernih robotskih ehografskih sustava. Admitancijska kontrola se ističe kao praktičan i stabilan izbor za kolaborativne robote s pozicijskim sučeljem, dok hibridni pristup omogućuje jasnu podjelu prostora zadatka na tangencijalno gibanje sonde i aksijalnu regulaciju sile.

U okviru naprednih pristupa (Varijanta C), analizirana je važnost vizualno vođenog upravljanja (Image-guided control) koje koristi samu ultrazvučnu sliku u zatvorenoj petlji, čime se sila dinamički prilagođava akustičkom kontaktu. Primjena učenja (Reinforcement Learning) omogućuje prilagodbu složenim nelinearnim deformacijama tkiva, dok se problem robusnosti na respiratorno gibanje rješava prediktivnim Kalmanovim filtrima i vision-haptic fuzijom. Sigurnost pacijenta ostaje primarni tehnički imperativ koji se jamči višestrukim redundantnim softverskim i hardverskim modulima, ugrađenim u predloženu multimodalnu arhitekturu.

Budućnost robotskih ultrazvučnih sustava leži u daljnjoj integraciji umjetne inteligencije za automatsku dijagnostiku i 3D rekonstrukciju organa, čime će se ehografski pregledi učiniti potpuno standardiziranim, objektivnim i dostupnima na daljinu putem telerobotike, uz maksimalnu sigurnost i udobnost za pacijenta.
